% Source: https://artofproblemsolving.com/community/c3h1845492_interesting_geometry_problem
% Stable identity: aops-topic-1845492-q01
\textbf{$A$ and $C$ are on opposite sides of line $XY$.} If they were on the same side, then (with $XY$ as a horizontal base of length $5$) $A$ would be $9$ units and $C$ $3$ units above the line, at horizontal distance $5$ from each other, so $AC=\sqrt{5^2+(9-3)^2}=\sqrt{61}<9=AX$. But every point of the square is at most $AC$ (the diagonal) away from $A$, so $X$ could not be inside the square. Hence $A$ and $C$ are on opposite sides.

\textbf{Completing a right triangle.} Extend $AX$ beyond $X$ to the point $Z$ with $XZ=3$. Since $XZ$ and $YC$ are both perpendicular to $XY$, have length $3$ and point to the same side of $XY$, the quadrilateral $XYCZ$ is a rectangle. So $ZC=XY=5$, $ZC\parallel XY$, and $\angle AZC=90^\circ$. Also $AZ=AX+XZ=9+3=12$.

\begin{asy}
usepackage("amsmath");
size(6cm);
defaultpen(fontsize(10pt));
real s=13/sqrt(2);
pair A0=(0,-9),X0=(0,0),Y0=(5,0),C0=(5,3),Z0=(0,3);
real ang=45-degrees(atan2(12,5));
transform T=rotate(ang)*shift(-A0);
pair A=T*A0,X=T*X0,Y=T*Y0,C=T*C0,Z=T*Z0;
pair B=(s,0),D=(0,s);
draw(A--B--C--D--cycle,linewidth(0.9));
draw(A--X--Y--C,blue+linewidth(1));
draw(X--Z--C,dashed);
draw(A--C,red+dashed);
void ra(pair P,pair U,pair V){real e=0.55;pair u=e*unit(U-P),v=e*unit(V-P);draw(P+u--P+u+v--P+v,linewidth(0.6));}
ra(X,A,Y);ra(Y,X,C);ra(Z,X,C);
dot(X);dot(Y);dot(Z);
label("$A$",A,SW);label("$B$",B,SE);label("$C$",C,NE);label("$D$",D,NW);
label("$X$",X,dir(X-Z)+dir(X-Y),blue);label("$Y$",Y,dir(Y-C)+dir(Y-X),blue);label("$Z$",Z,dir(Z-X)+dir(Z-C));
label("$9$",(A+X)/2,dir(rotate(90)*(X-A)),blue);label("$5$",(X+Y)/2,dir(rotate(-90)*(Y-X)),blue);label("$3$",Y+0.35*(C-Y),dir(rotate(-90)*(C-Y)),blue);
label("$3$",(X+Z)/2,dir(rotate(90)*(Z-X)));label("$5$",(Z+C)/2,dir(rotate(90)*(C-Z)));
label("\mbox{}",(min(currentpicture,user=true)+max(currentpicture,user=true))/2,black); // keeps the text after the figure black when LaTeX typesets the labels inline
\end{asy}

By the Pythagorean theorem in triangle $AZC$,
\[
AC^2=AZ^2+ZC^2=12^2+5^2=169 ,
\]
so the diagonal is $AC=13$. A square with diagonal $13$ has side $\dfrac{13}{\sqrt2}$ and area
\[
\dfrac{AC^2}{2}=\boxed{\dfrac{169}{2}} .
\]
(With these values $X$ and $Y$ do lie inside the square, as the figure shows.)
